% Ampliación estadística para el artículo V (ES).
% Fuente principal: N19; antecedente de memoria ordenada: F044.
% Este archivo no modifica la fuente activa y debe integrarse como sucesor.
\section{Dualidad de Fourier, incertidumbre finita y memoria ordenada}
\label{v:n19:sec:incertidumbre-fourier}

El dominio de la realización finita procede de la cadena tipada ya
construida en el artículo:
\begin{equation}
 \mathsf{APP}\longrightarrow\mathsf{TRIT}\longrightarrow\mathsf{TPK}
 \longrightarrow X_{\mathrm{enr}}
 \longrightarrow (\mathbb Z/9\mathbb Z)^3
 \xrightarrow{\,\beta_{729}\,}
 \Lambda:=(\mathbb Z/27\mathbb Z)^2
 \longrightarrow \ell^2(\Lambda).
 \label{v:n19:eq:cadena-tipificada}
\end{equation}
Aquí \(X_{\mathrm{enr}}\) conserva ruta, orientación, acarreo, frontera y
memoria; el paso a \((\mathbb Z/9\mathbb Z)^3\) es el registro posicional
pertinente. La aplicación \(\beta_{729}\), definida en
\eqref{v:ant:eq:ii-area-beta729} y demostrada biyectiva en la
Proposición~\ref{v:ant:prop:ii-area-beta729}, reagrupa los seis trits
ordenados en dos coordenadas de orden veintisiete. El registro posicional
y la realización de grupo quedan así diferenciados: \(\beta_{729}\)
transporta la incidencia ordenada, mientras que la transformada de
Fourier emplea las operaciones propias de
\((\mathbb Z/27\mathbb Z)^2\). La genealogía determina, por tanto, algo
más preciso que el cardinal común de \(729\) posiciones.

\subsection{Caracteres unitarios y bases mutuamente no sesgadas}
\label{v:n19:sec:caracteres-mub}

Escribamos \(N:=|\Lambda|=27^2=729\). Para
\(x=(x_1,x_2)\) y \(k=(k_1,k_2)\) en \(\Lambda\), sea
\begin{equation}
 \chi_k(x):=
 \exp\!\left(\frac{2\pi i}{27}(k_1x_1+k_2x_2)\right).
 \label{v:n19:eq:caracteres}
\end{equation}
El producto escalar de las coordenadas se toma módulo veintisiete. La
ortogonalidad de caracteres da
\begin{equation}
 \sum_{x\in\Lambda}\chi_k(x)\overline{\chi_{\ell}(x)}
 =N\,\delta_{k,\ell}.
 \label{v:n19:eq:ortogonalidad-caracteres}
\end{equation}
Definimos la transformada unitaria \(\mathcal F_\Lambda\) por
\begin{equation}
 \widehat\psi(k):=(\mathcal F_\Lambda\psi)(k)
 =\frac1{\sqrt N}\sum_{x\in\Lambda}\chi_k(x)\psi(x),
 \qquad
 \psi(x)=\frac1{\sqrt N}\sum_{k\in\Lambda}
 \overline{\chi_k(x)}\widehat\psi(k).
 \label{v:n19:eq:fourier-inversion}
\end{equation}
El carácter conjugado de la segunda fórmula es indispensable bajo la
convención positiva de la primera. La identidad
\eqref{v:n19:eq:ortogonalidad-caracteres} prueba simultáneamente la
inversión y \(\|\mathcal F_\Lambda\psi\|_2=\|\psi\|_2\).

\begin{lemma}[No sesgo mutuo de las bases]
\label{v:n19:lem:mub}
Sean \(\{e_x:x\in\Lambda\}\) la base de posición y
\(f_k(x)=N^{-1/2}\overline{\chi_k(x)}\) la base de Fourier. Entonces
\begin{equation}
 |\langle e_x,f_k\rangle|=N^{-1/2}=1/27
 \qquad(x,k\in\Lambda).
 \label{v:n19:eq:mub}
\end{equation}
Por tanto, ambas bases son mutuamente no sesgadas.
\end{lemma}
\begin{proof}
La ortonormalidad de \(\{f_k\}\) se sigue de
\eqref{v:n19:eq:ortogonalidad-caracteres}; además,
\(\langle e_x,f_k\rangle=f_k(x)\) y todo carácter tiene módulo uno.
\end{proof}

\subsection{Incertidumbre exacta de soporte}
\label{v:n19:sec:soporte}

Para una función \(u:\Lambda\to\mathbb C\), denotemos por
\(\operatorname{supp}u=\{x:u(x)\ne0\}\) su soporte.

\begin{theorem}[Cota de soporte en el toro triádico]
\label{v:n19:thm:soporte}
Para todo \(\psi\in\ell^2(\Lambda)\) no nulo,
\begin{equation}
 |\operatorname{supp}\psi|\,
 |\operatorname{supp}\widehat\psi|\ge N=729.
 \label{v:n19:eq:cota-soporte}
\end{equation}
\end{theorem}
\begin{proof}
Pongamos \(S=\operatorname{supp}\psi\) y
\(T=\operatorname{supp}\widehat\psi\). La inversión correcta de
\eqref{v:n19:eq:fourier-inversion} da, para todo \(x\in S\),
\[
 |\psi(x)|
 =\left|\frac1{\sqrt N}\sum_{k\in T}
       \overline{\chi_k(x)}\widehat\psi(k)\right|
 \le \sqrt{\frac{|T|}{N}}\,\|\widehat\psi\|_2,
\]
por Cauchy--Schwarz y \(|\chi_k(x)|=1\). Al elevar al cuadrado y
sumar sobre \(S\),
\[
 \|\psi\|_2^2
 \le\frac{|S||T|}{N}\,\|\widehat\psi\|_2^2
 =\frac{|S||T|}{N}\,\|\psi\|_2^2.
\]
La norma es positiva; su cancelación prueba
\eqref{v:n19:eq:cota-soporte}. Esta es la especialización al grupo
finito realizado de la incertidumbre de soporte de
Donoho--Stark~\cite{donoho_stark_1989}.
\end{proof}

La cadena \eqref{v:n19:eq:cadena-tipificada} fija el dominio antes de
formular la desigualdad: APP--TRIT--TPK aporta el registro posicional y
\(\beta_{729}\) su realización finita como grupo \(\Lambda\).

\subsection{Anuladores y clasificación de los casos de saturación}
\label{v:n19:sec:anuladores}

Si \(H\le\Lambda\), su anulador es
\begin{equation}
 H^\perp:=\{k\in\Lambda:\chi_k(h)=1\text{ para todo }h\in H\}.
 \label{v:n19:eq:anulador}
\end{equation}
Para divisores \(a,b\mid27\), consideremos
\begin{equation}
 H_{a,b}:=\{(x_1,x_2)\in\Lambda:
 x_1\equiv0\pmod a,\ x_2\equiv0\pmod b\}.
 \label{v:n19:eq:subgrupo-rectangular}
\end{equation}
Entonces
\begin{equation}
 |H_{a,b}|=\frac{27}{a}\frac{27}{b},
 \quad
 H_{a,b}^{\perp}
 =\{(k_1,k_2):k_1\equiv0\pmod{27/a},
                  k_2\equiv0\pmod{27/b}\},
 \quad |H_{a,b}^{\perp}|=ab.
 \label{v:n19:eq:anulador-explicito}
\end{equation}
En efecto, la suma geométrica en cada coordenada es nula salvo cuando
\(27\mid ak_1\) y \(27\mid bk_2\). Por consiguiente, para el indicador
normalizado
\begin{equation}
 \psi_{a,b}:=|H_{a,b}|^{-1/2}\mathbf1_{H_{a,b}},
 \label{v:n19:eq:indicador-normalizado}
\end{equation}
se obtiene exactamente
\begin{equation}
 \widehat\psi_{a,b}(k)
 =\sqrt{\frac{|H_{a,b}|}{N}}\,
   \mathbf1_{H_{a,b}^{\perp}}(k).
 \label{v:n19:eq:fourier-indicador}
\end{equation}
Así, la transformada es uniforme sobre el anulador y
\(|H_{a,b}|\,|H_{a,b}^{\perp}|=N\).

\begin{table}[htbp]
\centering
\caption{Cuatro saturaciones exactas, con dominios diferenciados.}
\label{v:n19:tab:saturaciones}
\begin{tabular}{@{}ccccc@{}}
\toprule
\((a,b)\) & \(|H_{a,b}|\) & \(|H_{a,b}^{\perp}|\) & producto
 & \((H_{\rm pos},H_{\rm F})\) \\
\midrule
\((3,3)\)  & \(81\) & \(9\)  & \(729\) & \((\log81,\log9)\) \\
\((9,9)\)  & \(9\)  & \(81\) & \(729\) & \((\log9,\log81)\) \\
\((3,9)\)  & \(27\) & \(27\) & \(729\) & \((\log27,\log27)\) \\
\((1,27)\) & \(27\) & \(27\) & \(729\) & \((\log27,\log27)\) \\
\bottomrule
\end{tabular}
\end{table}
Los dos últimos renglones poseen los mismos cardinales, pero no el
mismo subgrupo ni el mismo anulador; la igualdad numérica no colapsa sus
incidencias.

\subsection{Incertidumbre entrópica por interpolación y diferenciación}
\label{v:n19:sec:entropia}

Sea \(\|\psi\|_2=1\), y definamos las distribuciones
\begin{equation}
 p_x:=|\psi(x)|^2,\qquad q_k:=|\widehat\psi(k)|^2,
 \qquad
 H_{\rm pos}:=-\sum_xp_x\log p_x,\quad
 H_{\rm F}:=-\sum_kq_k\log q_k,
 \label{v:n19:eq:entropias-distribucion}
\end{equation}
con logaritmo natural y convenio \(0\log0=0\). Son entropías de
Shannon de dos distribuciones de medida sobre bases conjugadas.

\begin{theorem}[Cota entrópica de la realización \(\Lambda\)]
\label{v:n19:thm:entropia}
Todo estado unitario satisface
\begin{equation}
 H_{\rm pos}(\psi)+H_{\rm F}(\psi)
 \ge\log N=\log729=6\log3.
 \label{v:n19:eq:cota-entropica}
\end{equation}
\end{theorem}
\begin{proof}
De \(\|\mathcal F_\Lambda\|_{\ell^1\to\ell^\infty}
\le N^{-1/2}\) y la unitariedad
\(\|\mathcal F_\Lambda\|_{\ell^2\to\ell^2}=1\), el teorema de
interpolación de Riesz--Thorin~\cite{riesz_1927,thorin_1948} da, para
\(1<p\le2\) y \(p'=p/(p-1)\),
\begin{equation}
 \|\mathcal F_\Lambda\psi\|_{p'}
 \le N^{\frac12-\frac1p}\|\psi\|_p.
 \label{v:n19:eq:riesz-thorin}
\end{equation}
Consideremos
\begin{equation}
 g(p):=\log\|\mathcal F_\Lambda\psi\|_{p'}
       -\log\|\psi\|_p
       -\left(\frac12-\frac1p\right)\log N.
 \label{v:n19:eq:funcion-interpolacion}
\end{equation}
La desigualdad \eqref{v:n19:eq:riesz-thorin} implica \(g(p)\le0\),
mientras que \(g(2)=0\); por ello \(g'_-(2)\ge0\).

Para todo vector unitario \(u\) en dimensión finita,
\begin{equation}
 \left.\frac{d}{dr}\log\|u\|_r\right|_{r=2}
 =\frac14\sum_j|u_j|^2\log|u_j|^2
 =-\frac14H(|u|^2).
 \label{v:n19:eq:derivada-norma}
\end{equation}
Los términos nulos se obtienen por continuidad; los términos positivos
son diferenciables. Como \(\left.dp'/dp\right|_{p=2}=-1\), al derivar
\eqref{v:n19:eq:funcion-interpolacion} por la izquierda se obtiene
\begin{equation}
 g'_-(2)=\frac14
 \bigl(H_{\rm pos}+H_{\rm F}-\log N\bigr)\ge0,
 \label{v:n19:eq:derivada-entropica}
\end{equation}
que equivale a \eqref{v:n19:eq:cota-entropica}.

Esta derivación pertenece a la línea de incertidumbre entrópica de
Hirschman y Beckner~\cite{hirschman_1957,beckner_1975}; para dos bases
arbitrarias, la formulación por solapamiento máximo es la de
Maassen--Uffink~\cite{maassen_uffink_1988}. En la presente realización,
la constante \(\log N\) se obtiene directamente al diferenciar el factor
de interpolación de \eqref{v:n19:eq:riesz-thorin}.
\end{proof}

Para los estados \(\psi_{a,b}\), \eqref{v:n19:eq:fourier-indicador}
muestra que \(p\) y \(q\) son uniformes sobre \(H_{a,b}\) y
\(H_{a,b}^{\perp}\), respectivamente. Por tanto,
\[
 H_{\rm pos}=\log|H_{a,b}|,\qquad
 H_{\rm F}=\log|H_{a,b}^{\perp}|,\qquad
 H_{\rm pos}+H_{\rm F}=\log729,
\]
lo que prueba la saturación entrópica de los cuatro casos de la
Tabla~\ref{v:n19:tab:saturaciones} sin aproximación numérica.

\subsection{Dominios tipados de la información}
\label{v:n19:sec:tipos-entropia}

La cantidad \(H_{\rm pos}+H_{\rm F}\) mide la suma de dos entropías
clásicas de distribución asociadas a medidas conjugadas. La entropía de
entrelazamiento del corte de borde, la entropía de von Neumann de un
estado reducido, la entropía termodinámica y la temperatura, y la entropía
marginal del histograma trítico se definen sobre objetos y dominios
distintos. Toda comparación entre estas magnitudes se formula mediante un
mapa explícito que especifica dominio, codominio y estado. La cota
\eqref{v:n19:eq:cota-entropica} pertenece íntegramente a
\(\ell^2(\Lambda)\).

\subsection{Memoria ordenada en el dominio dodecafásico}
\label{v:n19:sec:memoria-ordenada}

El resultado F044 actúa sobre palabras dodecafásicas y caracteriza la
información de orden conservada por el estado enriquecido. Su dominio y
sus operadores son independientes de la realización de Fourier de N19;
la relación entre ambos resultados compara la información marginal con la
información ordenada.

Sea \(\tau=(\tau_1,\ldots,\tau_{12})\in\{-1,0,+1\}^{12}\), con
\(\operatorname{sgn}(0)=0\), e índices cíclicos módulo doce. Definamos
\begin{align}
 G_a&:=\{a,a+4,a+8\},\qquad a=1,2,3,4,\nonumber\\
 \nu_{120}(\tau)&:=\sum_{a=1}^4
   \operatorname{sgn}\!\left(\sum_{m\in G_a}\tau_m\right),
 \label{v:n19:eq:nu120}\\
 \nu_{270}(\tau)&:=\sum_{m=1}^{12}\tau_m\tau_{m+3}.
 \label{v:n19:eq:nu270}
\end{align}
Las palabras
\begin{equation}
 \tau^{(a)}=(1,1,0,0,0,0,0,0,0,0,0,0),\qquad
 \tau^{(b)}=(1,0,0,1,0,0,0,0,0,0,0,0)
 \label{v:n19:eq:testigos-orden}
\end{equation}
tienen el mismo histograma ---diez ceros y dos unidades--- y, por tanto,
la misma entropía marginal de alfabeto. No obstante,
\begin{equation}
 (\nu_{120},\nu_{270})(\tau^{(a)})=(2,0),\qquad
 (\nu_{120},\nu_{270})(\tau^{(b)})=(2,1).
 \label{v:n19:eq:testigos-lectura}
\end{equation}
El segundo lector detecta una incidencia a distancia tres que el
histograma elimina.

Para la involución \(C_M\tau:=-\operatorname{rev}(\tau)\), la reversión
permuta las cuatro clases \(G_a\) y la negación invierte sus sumas;
la autocorrelación cíclica a paso \(+3\) coincide, tras reindexar, con
la de paso \(-3\). En consecuencia,
\begin{equation}
 \nu_{120}(C_M\tau)=-\nu_{120}(\tau),\qquad
 \nu_{270}(C_M\tau)=\nu_{270}(\tau).
 \label{v:n19:eq:paridades-lectores}
\end{equation}
El lector \(\nu_{120}\) es impar y el lector cuadrático \(\nu_{270}\) es
par. Estas paridades caracterizan la información adicional que conserva
la palabra ordenada.

\subsection{Dominio de validez y condiciones matemáticas}
\label{v:n19:sec:falsadores}

Las condiciones que delimitan los dos resultados son:
\begin{enumerate}
 \item \(\beta_{729}\) es biyectiva sobre la realización posicional de
 \(\Lambda\) usada en \eqref{v:n19:eq:cadena-tipificada};
 \item la ortogonalidad \eqref{v:n19:eq:ortogonalidad-caracteres} y la
 inversión con carácter conjugado establecen la unitariedad;
 \item todo estado no nulo satisface el producto de soportes
 \(\lvert\operatorname{supp}\psi\rvert
 \lvert\operatorname{supp}\widehat\psi\rvert\ge729\);
 \item todo estado unitario satisface
 \(H_{\rm pos}+H_{\rm F}\ge\log729\);
 \item en los cuatro casos de la
 Tabla~\ref{v:n19:tab:saturaciones}, la transformada es uniforme sobre el
 anulador y el producto de cardinales es \(729\);
 \item los testigos de \eqref{v:n19:eq:testigos-orden} comparten
 histograma, producen \eqref{v:n19:eq:testigos-lectura} y satisfacen las
 paridades \eqref{v:n19:eq:paridades-lectores}.
\end{enumerate}
Los cinco primeros enunciados pertenecen a la realización finita tipada
\(\Lambda\) y a su Fourier unitaria; el sexto pertenece al dominio
dodecafásico de memoria ordenada. Estas condiciones proporcionan también
falsadores locales: el incumplimiento de una identidad, desigualdad o
valor testigo refuta el enunciado correspondiente en su dominio. La
relación entre N19 y F044 es complementaria y mantiene separados sus
dominios y operadores.
